% =========================================================================
% PREÁMBULO PARA ARTÍCULO CIENTÍFICO (FORMATO SPRINGER LNCS / SIMBIG)
% =========================================================================

\usepackage[T1]{fontenc}        % Codificación de fuentes de salida
\usepackage{lmodern}            % Fuentes escalables Latin Modern (requerido por microtype)
\usepackage[utf8]{inputenc}     % Codificación de caracteres UTF-8
\usepackage[spanish,es-tabla]{babel} % Idioma español con denominación "Tabla"
\AtBeginDocument{\shorthandoff{<>}}  % Evita interferencias con caracteres especiales

\usepackage{graphicx}           % Inclusión y manipulación de figuras
\usepackage{booktabs}           % Reglas tipográficas para tablas (toprule, midrule, bottomrule)
\usepackage{tabularx}           % Tablas con ajuste proporcional automático
\usepackage{array}              % Modificadores avanzados de columnas
\usepackage{enumerate}          % Listas con etiquetas personalizadas [(a)], [1.]
\usepackage{amsmath, amssymb}   % Símbolos y entornos matemáticos
\usepackage{microtype}          % Ajustes microtipográficos
\usepackage{url}                % Manejo de hipervínculos URL
\usepackage{xurl}               % División inteligente de URLs largas
\usepackage[table]{xcolor}      % Soporte de color en celdas de tablas

% Configuración de hipervínculos para artículo científico
\usepackage{hyperref}
\hypersetup{
    colorlinks=true,
    linkcolor=blue,
    citecolor=blue,
    filecolor=blue,
    urlcolor=blue,
    pdfstartview=FitH
}
